% =====================================================================
% 06b_partie_d_triggers.tex — triggers, transactions et clôture
% =====================================================================
\chapter{Partie D : triggers, transactions et clôture}
\label{chap:partieD2}

\section{Objectifs du chapitre}

Ce chapitre couvre la partie la plus sensible du projet : les triggers, les
transactions et la clôture de compte. Les mécanismes PL/SQL utilisés suivent
la référence officielle \cite{oraclePLSQL}. Ces mécanismes modifient l'état de la
base ; ils doivent donc être présentés dans un ordre contrôlé. Les triggers
assurent la sécurité des opérations, le trigger \code{AFTER} met à jour le
solde, \code{passeMvt} vérifie l'idempotence du recalcul, \code{virement}
garantit l'atomicité, et \code{fermeCompte} détruit des comptes. Cette dernière
procédure est volontairement testée en dernier.

\begin{attentionbox}[Chronologie]
Les tests \code{passeMvt} et \code{virement} précèdent \code{fermeCompte}. Les
tests de \code{fermeCompte} suppriment CE4 et CC1 ; ils ne doivent donc pas
être exécutés avant les démonstrations sur les intérêts, les opérations ou les
alertes. Cette contrainte expliquera plusieurs résultats apparemment
contradictoires entre la partie C et la fin de la partie D.
\end{attentionbox}

\section{Les six triggers}

Le projet comporte six triggers : un trigger \code{COMPOUND} de validation, un
trigger \code{AFTER INSERT} de mise à jour et d'alerte, un trigger d'échéance
et trois triggers d'audit. Le tableau~\ref{tab:triggers-D2} rappelle leur
événement et leur rôle.

\begin{table}[htbp]
  \centering
  \small
  \caption{Triggers de la partie D et responsabilités.}
  \label{tab:triggers-D2}
  \begin{tabularx}{\textwidth}{@{}L{4.6cm}L{4.9cm}X@{}}
    \toprule
    \entete{Trigger} & \entete{Événement} & \entete{Rôle} \\
    \midrule
    \code{trg\_op\_avant} &
    \code{BEFORE INSERT} sur \code{operation}, \code{COMPOUND} &
    Vérifie le statut, les droits, le plafond journalier et le découvert. \\
    \code{trg\_op\_apres} &
    \code{AFTER INSERT} sur \code{operation} &
    Met à jour le solde et \code{nb\_op\_cb}, crée les alertes. \\
    \code{trg\_ct\_echeance} &
    \code{AFTER INSERT OR UPDATE OF date\_echeance} sur \code{cpt\_terme} &
    Crée une alerte \code{ECHEANCE\_PROCHE} ou \code{ECHEANCE\_DEPASSEE}. \\
    \code{trg\_audit\_taux} &
    \code{AFTER UPDATE OF taux\_interet} sur \code{cpt\_epargne} &
    Trace l'ancien et le nouveau taux dans \code{"AUDIT"}. \\
    \code{trg\_audit\_taux\_terme} &
    \code{AFTER UPDATE OF taux\_garanti} sur \code{cpt\_terme} &
    Trace l'ancien et le nouveau taux garanti. \\
    \code{trg\_audit\_decouvert} &
    \code{AFTER UPDATE OF decouvert} sur \code{cpt\_courant} &
    Trace l'ancien et le nouveau découvert. \\
    \bottomrule
  \end{tabularx}
\end{table}

\section{Le trigger \code{COMPOUND}}

Le trigger \code{trg\_op\_avant} doit vérifier quatre règles : le compte est
actif, le client a le droit de débiter, le plafond journalier de 200\,000~DA
n'est pas dépassé et le découvert autorisé n'est pas franchi. La difficulté
vient du plafond : il faut additionner les retraits du jour, donc lire la table
\code{operation} pendant un insert sur \code{operation}. Un trigger ligne par
ligne produirait \code{ORA-04091}, l'erreur de table mutante.

La solution est un trigger \code{COMPOUND}. La section
\code{BEFORE STATEMENT} prend un instantané des retraits du jour dans une
collection associative indexée par numéro de compte. La section
\code{BEFORE EACH ROW} vérifie les règles et accumule le montant de la ligne
courante. Elle ne relit jamais \code{operation}.

\begin{lstlisting}[caption={Trigger \code{COMPOUND} de validation des opérations.},label={lst:trg-avant}]
CREATE OR REPLACE TRIGGER trg_op_avant
FOR INSERT ON operation
COMPOUND TRIGGER
  c_plafond CONSTANT NUMBER := 200000;
  TYPE t_retraits IS TABLE OF NUMBER INDEX BY VARCHAR2(10);
  v_retraits t_retraits;
  v_num      VARCHAR2(10);
  v_statut   VARCHAR2(10);
  v_solde    NUMBER;
  v_decouvert NUMBER;
  v_nb       NUMBER;
  v_cumul    NUMBER;

  BEFORE STATEMENT IS
  BEGIN
    FOR r IN (SELECT cc.num_compte, SUM(-o.montant) AS total
                FROM operation o, cpt_courant cc
               WHERE o.ref_compte = REF(cc)
                 AND o.montant < 0
                 AND o.date_op >= TRUNC(SYSDATE)
               GROUP BY cc.num_compte) LOOP
      v_retraits(r.num_compte) := r.total;
    END LOOP;
  END BEFORE STATEMENT;

  BEFORE EACH ROW IS
  BEGIN
    SELECT x.num_compte, x.statut, x.solde, x.decouvert
      INTO v_num, v_statut, v_solde, v_decouvert
      FROM cpt_courant x
     WHERE REF(x) = :NEW.ref_compte;

    IF v_statut <> 'ACTIF' THEN
      RAISE_APPLICATION_ERROR(-20001, 'Compte non actif.');
    END IF;

    IF :NEW.montant < 0 THEN
      SELECT COUNT(*) INTO v_nb
        FROM cpt_courant x
       WHERE REF(x) = :NEW.ref_compte
         AND x.ref_client = :NEW.ref_client;

      IF v_nb = 0 THEN
        SELECT COUNT(*) INTO v_nb
          FROM cpt_courant x, TABLE(x.signataires) s
         WHERE REF(x) = :NEW.ref_compte
           AND s.ref_client = :NEW.ref_client
           AND s.droit IN ('D', 'X');
        IF v_nb = 0 THEN
          RAISE_APPLICATION_ERROR(-20002, 'Client sans droit de debit.');
        END IF;
      END IF;

      v_cumul := CASE WHEN v_retraits.EXISTS(v_num)
                      THEN v_retraits(v_num) ELSE 0 END;
      IF v_cumul - :NEW.montant > c_plafond THEN
        RAISE_APPLICATION_ERROR(-20003, 'Plafond journalier depasse.');
      END IF;
      v_retraits(v_num) := v_cumul - :NEW.montant;

      IF v_solde + :NEW.montant - :NEW.commission
         < -NVL(v_decouvert, 0) THEN
        RAISE_APPLICATION_ERROR(-20004, 'Decouvert autorise depasse.');
      END IF;
    END IF;
  END BEFORE EACH ROW;
END;
/
\end{lstlisting}

\begin{clebox}[Les quatre codes d'erreur]
\code{ORA-20001} signale un compte non actif ; \code{ORA-20002} un défaut de
droit ; \code{ORA-20003} un dépassement du plafond journalier ; \code{ORA-20004}
un découvert dépassé. Ces codes sont applicatifs et volontairement distincts
des codes Oracle standard. Ils rendent les tests d'exception lisibles et
permettent de vérifier que le bon contrôle a échoué.
\end{clebox}

\section{Le trigger \code{AFTER INSERT}}

Le trigger \code{trg\_op\_apres} applique le mouvement au solde :
\code{solde = solde + montant - commission}. Il incrémente \code{nb\_op\_cb}
lorsque le type vaut \code{CB}. Il crée une alerte \code{SOLDE\_NEGATIF} si le
solde passe de positif ou nul à négatif, et une alerte
\code{OPERATION\_SUSPECTE} si le montant dépasse 1\,000\,000~DA.

\begin{lstlisting}[caption={Trigger \code{AFTER INSERT} : solde, compteur CB et alertes.},label={lst:trg-apres}]
CREATE OR REPLACE TRIGGER trg_op_apres
AFTER INSERT ON operation
FOR EACH ROW
DECLARE
  v_num         VARCHAR2(10);
  v_solde_avant NUMBER;
  v_solde_apres NUMBER;
  v_nom         VARCHAR2(30);
BEGIN
  SELECT x.num_compte, x.solde
    INTO v_num, v_solde_avant
    FROM cpt_courant x
   WHERE REF(x) = :NEW.ref_compte;

  v_solde_apres := v_solde_avant + :NEW.montant - :NEW.commission;

  UPDATE cpt_courant
     SET solde = v_solde_apres,
         nb_op_cb = NVL(nb_op_cb, 0)
                    + CASE WHEN :NEW.type_op = 'CB' THEN 1 ELSE 0 END
   WHERE num_compte = v_num;

  SELECT c.nom INTO v_nom
    FROM client c
   WHERE REF(c) = :NEW.ref_client;

  IF v_solde_avant >= 0 AND v_solde_apres < 0 THEN
    INSERT INTO alerte (id_alerte, type_alerte, num_compte,
                        message, date_alerte, traitee)
    VALUES (seq_alerte.NEXTVAL, 'SOLDE_NEGATIF', v_num,
            'Le compte '||v_num||' de '||v_nom
            ||' passe en solde negatif ('||v_solde_apres||' DA).',
            SYSDATE, 'N');
  END IF;

  IF :NEW.montant > 1000000 THEN
    INSERT INTO alerte (id_alerte, type_alerte, num_compte,
                        message, date_alerte, traitee)
    VALUES (seq_alerte.NEXTVAL, 'OPERATION_SUSPECTE', v_num,
            'Operation de '||:NEW.montant||' DA (> 1000000) par '
            ||v_nom||'.', SYSDATE, 'N');
  END IF;
END;
/
\end{lstlisting}

\begin{remarquebox}[Pourquoi le solde est-il mis à jour dans un trigger ?]
Le trigger garantit que toute insertion dans \code{operation} met à jour le
solde, y compris une insertion SQL directe qui ne passerait pas par
\code{passeMvt}. La méthode \code{passeMvt} recalcule ensuite le solde depuis
l'historique ; comme le trigger a déjà appliqué le mouvement, le recalcul
redonne le même résultat. Il est idempotent et ne double pas le montant.
\end{remarquebox}

\section{Le trigger d'échéance}

Le trigger \code{trg\_ct\_echeance} se déclenche après insertion ou modification
de \code{date\_echeance} sur \code{cpt\_terme}. Si l'échéance est dépassée, il
crée une alerte \code{ECHEANCE\_DEPASSEE} ; si elle est dans huit jours ou
moins, il crée une alerte \code{ECHEANCE\_PROCHE}. Un test
\code{COUNT(*) = 0} sur \code{alerte}, pour le même compte, le même type et
\code{traitee = 'N'}, évite les doublons.

\begin{lstlisting}[caption={Trigger d'alerte sur l'échéance d'un compte à terme.},label={lst:trg-echeance}]
CREATE OR REPLACE TRIGGER trg_ct_echeance
AFTER INSERT OR UPDATE OF date_echeance ON cpt_terme
FOR EACH ROW
DECLARE
  v_jours NUMBER;
  v_nb    NUMBER;
BEGIN
  v_jours := TRUNC(:NEW.date_echeance) - TRUNC(SYSDATE);

  IF v_jours < 0 THEN
    SELECT COUNT(*) INTO v_nb
      FROM alerte
     WHERE num_compte = :NEW.num_compte
       AND type_alerte = 'ECHEANCE_DEPASSEE'
       AND traitee = 'N';
    IF v_nb = 0 THEN
      INSERT INTO alerte (id_alerte, type_alerte, num_compte,
                          message, date_alerte, traitee)
      VALUES (seq_alerte.NEXTVAL, 'ECHEANCE_DEPASSEE', :NEW.num_compte,
              'Echeance depassee de '||(-v_jours)||' jour(s).',
              SYSDATE, 'N');
    END IF;
  ELSIF v_jours <= 8 THEN
    SELECT COUNT(*) INTO v_nb
      FROM alerte
     WHERE num_compte = :NEW.num_compte
       AND type_alerte = 'ECHEANCE_PROCHE'
       AND traitee = 'N';
    IF v_nb = 0 THEN
      INSERT INTO alerte (id_alerte, type_alerte, num_compte,
                          message, date_alerte, traitee)
      VALUES (seq_alerte.NEXTVAL, 'ECHEANCE_PROCHE', :NEW.num_compte,
              'Echeance dans '||v_jours||' jour(s).', SYSDATE, 'N');
    END IF;
  END IF;
END;
/
\end{lstlisting}

\section{Les triggers d'audit}

Les trois triggers d'audit suivent le même modèle : ils insèrent dans
\code{"AUDIT"} le nom de la table, le nom de la colonne, le numéro de compte,
l'ancienne valeur, la nouvelle valeur, la date et l'utilisateur. La table
\code{"AUDIT"} doit être écrite entre guillemets, car \code{AUDIT} est un mot
réservé.

\begin{lstlisting}[caption={Audit du découvert autorisé.},label={lst:trg-audit}]
CREATE OR REPLACE TRIGGER trg_audit_decouvert
AFTER UPDATE OF decouvert ON cpt_courant
FOR EACH ROW
BEGIN
  INSERT INTO "AUDIT" (id_audit, nom_table, nom_colonne, num_compte,
                       ancienne_valeur, nouvelle_valeur, date_modif,
                       utilisateur)
  VALUES (seq_audit.NEXTVAL, 'CPT_COURANT', 'DECOUVERT',
          :OLD.num_compte, :OLD.decouvert, :NEW.decouvert,
          SYSDATE, USER);
END;
/
\end{lstlisting}

\begin{attentionbox}[Audit et rollback]
Les triggers d'audit écrivent dans la même transaction que la modification.
Comme les tests 13 à 15 se terminent par \code{ROLLBACK}, les lignes d'audit
correspondantes disparaissent. À la fin du script, la table \code{"AUDIT"} est
donc vide, même si les triggers ont bien fonctionné. C'est un point à expliquer
clairement : l'audit est transactionnel, il n'est pas autonome.
\end{attentionbox}

\section{Résultats des tests D9}

Seize tests couvrent les triggers. Les cinq premiers vérifient des refus ; les
suivants vérifient des succès, des alertes et l'audit. Le
tableau~\ref{tab:tests-D9} reprend les résultats attendus.

\begin{table}[htbp]
  \centering
  \small
  \caption{Tests des triggers et résultats attendus.}
  \label{tab:tests-D9}
  \begin{tabularx}{\textwidth}{@{}L{6.6cm}L{4.4cm}X@{}}
    \toprule
    \entete{Test} & \entete{Résultat attendu} & \entete{Commentaire} \\
    \midrule
    1 : opération sur CC5 \code{BLOQUE} & \code{ORA-20001} &
    Compte non actif. \\
    2 : client 6 sur CC1, sans droit de débit & \code{ORA-20002} &
    Ni titulaire ni signataire autorisé. \\
    3 : client 5 (droit \code{R}) sur CC7 & \code{ORA-20002} &
    Droit de consultation seulement. \\
    4 : retrait de 60\,000 sur CC3 & \code{ORA-20004} &
    Découvert dépassé. \\
    5 : retrait de 210\,000 par le signataire \code{X} sur CC4 &
    \code{ORA-20003} & Plafond de 200\,000. \\
    6 : retrait de 1\,000 par le signataire \code{D} sur CC7 &
    Succès, 185\,000 $\rightarrow$ 184\,000 & Droit \code{D} accepté. \\
    7 : retrait de 1\,000 par le signataire \code{X} sur CC4 &
    Succès, 325\,000 $\rightarrow$ 324\,000 & Droit \code{X} accepté. \\
    8 : CB sur CC6 & 40\,000 $\rightarrow$ 39\,000 ; \code{nb\_op\_cb}
    0 $\rightarrow$ 1 & Compteur CB. \\
    9 : alerte \code{SOLDE\_NEGATIF} & 1 alerte pour CC6 &
    Transitoire, annulée par \code{ROLLBACK}. \\
    10 : alerte \code{OPERATION\_SUSPECTE} & 1 alerte pour CC1 &
    Crédit de plus de 1\,000\,000. \\
    11 : échéance proche CT99 & 1 alerte \code{ECHEANCE\_PROCHE} &
    Transitoire. \\
    12 : échéance dépassée CT98 & 1 alerte \code{ECHEANCE\_DEPASSEE} &
    Transitoire. \\
    13--15 : audit & Lignes dans \code{"AUDIT"} &
    Taux CE1, découvert CC1, taux garanti CT1 ; annulées par rollback. \\
    16 : \code{scan\_echeances} & Alertes CT2/CT3, sans doublon &
    Seules alertes persistantes. \\
    \bottomrule
  \end{tabularx}
\end{table}

\begin{attentionbox}[Ce qui reste après les tests]
Les tests 6 à 15 se terminent par \code{ROLLBACK}. Les alertes
\code{SOLDE\_NEGATIF}, \code{OPERATION\_SUSPECTE}, \code{ECHEANCE\_PROCHE}
(CT99) et \code{ECHEANCE\_DEPASSEE} (CT98) sont donc transitoires. Après le
script, il ne reste que les alertes CT2/CT3 créées par
\code{scan\_echeances}, et l'audit final est vide.
\end{attentionbox}

\section{\code{passeMvt} et l'idempotence}

La méthode \code{passeMvt} insère une opération, laisse le trigger
\code{AFTER} mettre à jour le solde, puis recalcule le solde à partir de
l'historique. Le recalcul est volontairement redondant : il garantit la
cohérence si la méthode est appelée après une insertion directe ou si le
trigger est désactivé. Il ne double pas le montant, car il remplace le solde par
la somme des mouvements et non par une addition.

\begin{lstlisting}[caption={Recalcul idempotent du solde dans \code{passeMvt}.},label={lst:recalcul-solde}]
UPDATE cpt_courant
   SET solde = (SELECT NVL(SUM(o.montant - o.commission), 0)
                  FROM operation o
                 WHERE o.ref_compte = REF(cpt_courant))
 WHERE num_compte = v_num;
\end{lstlisting}

Les tests sont exécutés dans des blocs séparés, avec un \code{ROLLBACK} entre
chaque essai. CC7 part de 185\,000. Un crédit de 10\,000 donne 195\,000 ; un
débit de 5\,000, exécuté dans un autre bloc, donne 180\,000. Un montant nul
lève \code{ORA-20030}, code distinct de \code{ORA-20010} utilisé par la
mutation de la vue objet.

\begin{clebox}[Pourquoi un code distinct pour le montant nul ?]
Le code \code{ORA-20010} est déjà utilisé par la méthode \code{mutation} pour
signaler qu'un conseiller n'est pas affecté à l'agence source. Utiliser le même
code pour un montant nul créerait une ambiguïté dans les tests. Le code
\code{ORA-20030} isole clairement le cas « montant nul interdit ».
\end{clebox}

\section{Le virement et l'atomicité}

La procédure \code{virement} a été présentée au chapitre précédent. Deux tests
complètent sa validation. Le premier, intra-agence, transfère 1\,000~DA de CC4
vers CC7 : la commission est nulle, CC4 passe de 325\,000 à 324\,000 et CC7 de
185\,000 à 186\,000. Le second, inter-agences, transfère 1\,000~DA de CC4 (A1)
vers CC3 (A3) : la commission est de 100~DA, CC4 passe de 325\,000 à 323\,900
et CC3 de 38\,000 à 39\,000.

Le test de rollback transfère 60\,000~DA de CC3 vers CC4. Le découvert de CC3
est insuffisant : le trigger lève \code{ORA-20004}, la procédure revient au
\code{SAVEPOINT} et CC4 reste à 325\,000. Ce test prouve que les deux écritures
sont annulées ensemble.

\begin{remarquebox}[Commission et trace comptable]
La commission inter-agences est portée par l'opération de débit, avec un
montant de 1~\% du virement et un minimum de 100~DA. Le crédit reste sans
commission. On évite ainsi une troisième opération de frais et on conserve la
trace du prélèvement sur le compte source.
\end{remarquebox}

\section{La clôture de compte}

\code{fermeCompte} est une procédure statique de \code{compte\_type}. Elle
cherche le compte dans \code{cpt\_courant}, puis dans \code{cpt\_epargne}, puis
dans \code{cpt\_terme}. Elle affiche le solde avant clôture, supprime les
opérations si c'est un compte courant, supprime le compte et valide par
\code{COMMIT}. Un compte inexistant déclenche \code{NO\_DATA\_FOUND} et un
message ; les autres erreurs provoquent un \code{ROLLBACK}.

\begin{lstlisting}[caption={Clôture d'un compte courant et suppression de ses opérations.},label={lst:ferme-compte}]
SELECT COUNT(*) INTO v_nb
  FROM cpt_courant WHERE num_compte = cpt;

IF v_nb > 0 THEN
  SELECT solde INTO v_solde
    FROM cpt_courant WHERE num_compte = cpt;
  DBMS_OUTPUT.PUT_LINE('fermeCompte : compte courant '||cpt
                       ||' - solde avant = '||v_solde);
  DELETE FROM operation
   WHERE ref_compte = (SELECT REF(x) FROM cpt_courant x
                        WHERE x.num_compte = cpt);
  DELETE FROM cpt_courant WHERE num_compte = cpt;
END IF;
\end{lstlisting}

Les tests destructifs sont exécutés en dernier. \code{fermeCompte('CE4')}
supprime un compte d'épargne de 320\,000~DA. \code{fermeCompte('CC1')} supprime
un compte courant de 2\,498\,500~DA et ses cinq opérations. Les appels sur
\code{CE9} et \code{XX1} vérifient le message « compte inexistant ». Après ces
tests, le client 4 n'a plus de compte d'épargne et le solde total des clients
change : c'est la raison pour laquelle les valeurs de la partie C ne doivent
pas être relues après la partie D.

\section{Chronologie et état final}

Le tableau~\ref{tab:chronologie-D} récapitule les disparitions d'objets et
leurs conséquences sur les résultats. Il constitue un aide-mémoire pour la
soutenance : chaque chiffre doit être associé à un moment précis.

\begin{table}[htbp]
  \centering
  \small
  \caption{Chronologie des tests destructifs et conséquences.}
  \label{tab:chronologie-D}
  \begin{tabularx}{\textwidth}{@{}L{3.8cm}L{5.2cm}X@{}}
    \toprule
    \entete{Moment} & \entete{Action} & \entete{Conséquence sur les résultats} \\
    \midrule
    Partie C & Requêtes 6.1 à 6.25 & CC1, CE4 et CT4 encore présents ;
    l'opération de 1\,500\,000 existe. \\
    Requête 6.25 & \code{DELETE} des CT expirés & CT4 supprimé ; 3 CT restants. \\
    Tests D9 & Triggers & Alertes et audits transitoires, annulés par
    \code{ROLLBACK}, sauf CT2/CT3. \\
    \code{passeMvt} & Crédit/débit sur CC7 & Tests séparés : 195\,000 puis
    180\,000, avec rollback. \\
    \code{virement} & CC4 $\rightarrow$ CC7, puis CC4 $\rightarrow$ CC3 &
    Soldes modifiés ; le test CC3 $\rightarrow$ CC4 échoue et ne modifie
    rien. \\
    \code{fermeCompte('CE4')} & Clôture de CE4 & CE4 supprimé ;
    \code{nbCepargne} du client 4 devient 0. \\
    \code{fermeCompte('CC1')} & Clôture de CC1 & CC1 et ses cinq opérations
    supprimés ; \code{soldeTotal} du client 1 diminue. \\
    \bottomrule
  \end{tabularx}
\end{table}

\begin{attentionbox}[Ne pas confondre état initial et état final]
Les valeurs de la partie C sont des valeurs \emph{avant} les tests destructifs.
Les valeurs finales sont : 6 clients, 6 comptes courants, 4 comptes d'épargne,
3 comptes à terme, 15 opérations, 9 signataires, 2 alertes et 9 affectations.
L'audit final est vide. Ces deux jeux de valeurs sont corrects ; ils
correspondent simplement à des moments différents.
\end{attentionbox}

\section{Pièges de la partie D2}

\begin{table}[htbp]
  \centering
  \small
  \caption{Pièges des triggers et des transactions.}
  \label{tab:pieges-D2}
  \begin{tabularx}{\textwidth}{@{}L{5.0cm}L{4.6cm}X@{}}
    \toprule
    \entete{Piège} & \entete{Symptôme} & \entete{Solution} \\
    \midrule
    Trigger ligne par ligne qui lit \code{operation} & \code{ORA-04091} &
    Trigger \code{COMPOUND}. \\
    \code{ALTER TYPE ADD MEMBER} pour un override &
    \code{ORA-22324} / \code{PLS-00636} &
    Déclarer \code{OVERRIDING} dès le \code{CREATE TYPE}. \\
    \code{ALTER TYPE} avec plusieurs membres & Erreur de syntaxe &
    Un membre par instruction. \\
    \code{ALTER TYPE} rejoué & \code{ORA-22324} &
    Protéger avec \code{SQLCODE = -22324}. \\
    Objets dépendants invalides après \code{CASCADE} &
    \code{STATUS = INVALID} & Recompiler les procédures et fonctions. \\
    \code{COMMIT} dans \code{virement} ou \code{passeMvt} &
    Le test de rollback échoue & Ne pas valider dans ces routines. \\
    \code{valeur(2)} sur un \code{VARRAY} de un élément &
    \code{ORA-06533} & Vérifier \code{COUNT} avant l'accès. \\
    \bottomrule
  \end{tabularx}
\end{table}

\section{Conclusion de la partie D2}

Les triggers transforment les règles métier en contrôles automatiques. Le
\code{COMPOUND} évite la table mutante, le trigger \code{AFTER} centralise la
mise à jour du solde et des alertes, et les triggers d'audit conservent la
trace des modifications sensibles. Les transactions sont pilotées par
l'appelant pour \code{virement} et \code{passeMvt}, mais autonomes pour
\code{drop\_client} et \code{fermeCompte}. Cette distinction est cohérente avec
la nature de chaque opération et constitue un point central de la soutenance.
